\section*{Author circuit conventions and problems}
A Boolean circuit consists of fan-in two AND gates, fan-in two OR gates, and NOT gates. The input gates are labelled by variables $x_1,\ldots,x_n$. We measure the size of a Boolean circuit $C$, denoted $|C|$, using the number of AND, OR, and NOT gates in the circuit. (While this is the convention adopted here, our techniques are robust to modifications of the circuit size measure and of the gate types in the circuit.)
The components of a multi-output Boolean function $f:\{0,1\}^n\to\{0,1\}^m$ are the single-output functions that compute the $i$th output bit of $f$ for $i\in[m]$. We say a Boolean circuit $C$ computes a multi-output Boolean function $f$ if for each component $f_i$ of $f$, there is a gate or input wire in $C$ that computes $f_i$.
\setcounter{nttheorem}{0}
\begin{ntdefinition}[Author Definition1]
Multi-MCSP is defined as follows:
\begin{description}
\item[Input.] Positive integers $n,m$, and $s$ represented in unary, and a (multi-output) Boolean function $f:\{0,1\}^n\to\{0,1\}^m$ represented by a string of length $m\cdot2^n$.
\item[Output.] The input is accepted if and only if there exists a Boolean circuit $C$ of size at most $s$ that computes $f$.
\end{description}
\end{ntdefinition}
Note that Multi-MCSP is in NP,\footnote{If the parameter $s$ is large then the answer is trivial.} and that this problem is at least as hard as MCSP.
\paragraph{Boolean functions over an arbitrary set of generators}
We also consider circuit minimization for (single-output, total) Boolean functions $f:\{0,1\}^m\to\{0,1\}$ under an arbitrary set of generators.
To explain, let $V$ be a finite set, called the ground set, and $\mathcal B=\{B_i\}_{i\in[m]}$ a family of nonempty sets $B_i\subseteq V$ called generators, and let $A\subseteq V$. Then we let $D(A\mid\mathcal B)$ denote the minimum number of unions, intersections, and complements that are required to construct $A$ from the sets in $\mathcal B$. More precisely, the complement of a set $U\subseteq V$ is defined as $V\setminus U$, and we represent a construction of $A$ from $\mathcal B$ as a sequence $B_1,\ldots,B_m,E_1,\ldots,E_\ell$ of sets in $V$ such that $E_\ell=A$ and each set $E_j$ is either the union or intersection of two previously generated sets, or the complement of a previously generated set. We assume for convenience that $D(B\mid\mathcal B)=0$ if $B\in\mathcal B$. We refer to $D(A\mid\mathcal B)$ as the discrete complexity\footnote{This general setting was already considered in~\cite{ref68}, see also Section1.7.2 of Jukna's book~\cite{ref42}. By Stone's representation theorem for Boolean algebras, discrete complexity can be seen as the investigation of circuit complexity with respect to an arbitrary Boolean algebra.} of $A$ with respect to $\mathcal B$.
It may then be seen the minimum number of AND, OR and NOT gates needed for a Boolean circuit to compute a Boolean function $f:\{0,1\}^m\to\{0,1\}$ is exactly $D(f^{-1}(1)\mid\mathcal B)$, where $V=\{0,1\}^m$, and $\mathcal B=\{x_1,\ldots,x_m\}\subseteq\{0,1\}^m$ contains the input variables $x_1\ldots,x_m$, seen as subsets of $\{0,1\}^m$ (i.e.. $x_i$ is the set of strings $w\in\{0,1\}^m$ such that $w_i=1$). So computing the discrete complexity $D(A\mid\mathcal B)$, for given $A$ and $\mathcal B$, generalizes the task of computing the minimum circuit size, by allowing for the consideration of generator sets $\mathcal B$ other than $\{x_1,\ldots,x_m\}$.
\setcounter{nttheorem}{2}
\begin{ntdefinition}[Author Definition3]
MDCP is defined as follows:
\begin{description}
\item[Input.] A positive integer $n$ represented in unary describing the size of the ground set $V=[n]$, a target set $A\subseteq V$, a family $\mathcal B$ of nonempty subsets of $V$, and an integer $s$.
\item[Output.] The input is accepted if and only if $D(A\mid\mathcal B)\leq s$.
\end{description}
\end{ntdefinition}
It is easy to see that MDCP is in NP and that it is more general than MCSP. Our hardness result for MDCP, discussed below, also holds under the assumption that the ground set $V$ is a hypercube $\{0,1\}^m$ and that the collection $\mathcal B$ contains the sets generated by $x_1,\ldots,x_m$.
\section*{2 Preliminaries}
We let $[n]$ denote the set $\{1,\ldots,n\}$. We will use the NP-hardness of approximating Set Cover with respect to sets of bounded size~\cite{ref76}. A weaker version of the result from~\cite{ref76} is sufficient.
We say that sets $S_1,\ldots,S_\ell$ cover a set $T$ if $T\subseteq S_1\cup\cdots\cup S_\ell$. For a collection of sets $\mathcal S$ and a set $T$, we use $\operatorname{cover}(T,\mathcal S)$ to denote the minimum number of sets in $\mathcal S$ necessary to cover $T$.
\setcounter{nttheorem}{6}
\begin{ntdefinition}[Author Definition7: $r$-Bounded Set Cover Problem]
For a positive integer $r$, the $r$-Bounded Set Cover Problem is defined as follows:
\begin{description}
\item[Input.] A positive integer $n$ represented in unary, and a collection $\mathcal S$ of nonempty subsets of $[n]$. We are promised that $\bigcup_{S\in\mathcal S}S=[n]$ and that $|S|\leq r$ for each $S\in\mathcal S$.
\item[Output.] The value $\operatorname{cover}([n],\mathcal S)$.
\end{description}
\end{ntdefinition}
For convenience, we defined the $r$-Bounded Set Cover Problem as an optimization problem instead of decision problem.
\begin{nttheorem}[Author Theorem8: Hardness of approximating $r$-Bounded Set Cover, {\cite{ref76}}]
For every constant $\alpha\geq1$ there exists $r\in\mathbb N$ such that approximating the $r$-Bounded Set Cover Problem within a factor of $\alpha$ is NP-hard. More precisely, for every $L\in\operatorname{NP}$, there exists a polynomial-time algorithm that, on input $x$, outputs a parameter $k$ and an instance $(1^m,\mathcal S)$ of the $r$-Bounded Set Cover Problem such that if $x\in L$ then $\operatorname{cover}([m],\mathcal S)\leq k$, and if $x\notin L$ then $\operatorname{cover}([m],\mathcal S)>\alpha\cdot k$.
\end{nttheorem}
